\section{The state of a spatial forestry market}\label{the-state-of-a-spatial-forestry-market}

\subsection{1. The central modelling problem}\label{the-central-modelling-problem}

A forestry market couples biological growth, harvesting, processing constraints,
transport, and information. Biological state determines feasible material.
Harvesting maps trees into logs. Processor specifications assign value to
regions of log space. Spatial operators determine delivered cost. Observations
update the decision maker's information.

A useful mathematical model must therefore preserve heterogeneity long enough
for the economically important compatibility tests to be applied.

Suppose two estates each contain 100 hectares and 20,000 cubic metres of standing
volume. Their commercial values can differ sharply. Estate A could contain mostly
logs above a target processor's small-end
diameter threshold and lie 20 km from a mill. Estate B could contain the same
total volume in smaller stems, more defect, and a 250 km transport route. The
state variables that explain the difference have been averaged away by the
scalar summary.

The paper's response is to make the primary state \textbf{measure-valued}. A measure
records how much mass or how many objects occupy each region of a state space.
For forestry, that state space can carry location, age, log-diameter and height,
species or genetics, quality, and management condition.

The standing-tree object is written schematically as

\[
\mu_t \in \mathcal M_+(E_T),
\]

where \texttt{E\_T} is the marked tree space and \texttt{M\_+(E\_T)} is the cone of finite
nonnegative Radon measures on that space.

The total mass \texttt{mu\_t(E\_T)} represents an intensity, count, or area-weighted count,
according to the sampling convention. Normalization would discard that quantity.

\subsection{2. Why there are three nested problems}\label{why-there-are-three-nested-problems}

\subsubsection{2.1 State estimation}\label{state-estimation}

Plots, remote-sensing features, harvest records, processor intake, and contractor
performance provide partial observations of the latent economic state.

The model therefore includes an observation process and a posterior state. A map
pixel supplies evidence about a latent tree distribution; a calibrated
observation model is required to infer DBH, log grade, or harvest value.

\subsubsection{2.2 Matching}\label{matching}

Once harvestable material has been represented as a log distribution, economic
usefulness is determined by compatibility. Processor \texttt{j} may accept only a
subset \texttt{S\_j} of log space. Operator \texttt{k} may have finite capacity or equipment
constraints. Delivered cost depends on the log, route, and operator.

An average-price identity such as

\[
\text{revenue} = \text{tonnes} \times \text{average price}.
\]

omits the relevant constraints. The allocation problem acts on feasible triples

\[
(\lambda, j, k),
\]

where \texttt{lambda} is a log state, \texttt{j} a processor, and \texttt{k} an operator.

\subsubsection{2.3 Sequential control}\label{sequential-control}

A decision maker can change both the current physical/economic outcome and the
quality of future information. Choosing an inventory campaign, instrumenting a
contractor, testing a processor's grade response, or exposing a new supply
cluster to the market can reveal parameters that affect later decisions.

This makes the problem a dual-control problem: actions can have both immediate
payoff and learning value.

\subsection{3. The operator chain}\label{the-operator-chain}

The paper is easier to understand when written as a sequence of operators.

\subsubsection{3.1 Biological evolution}\label{biological-evolution}

A stochastic evolution operator moves the standing-tree measure through size and
quality space and changes its mass through mortality and recruitment.

\subsubsection{3.2 Harvest and bucking}\label{harvest-and-bucking}

A harvest rule removes mass from the standing-tree measure. A bucking kernel
maps each harvested tree state to a finite measure of logs.

Schematically,

\[
\mu_t \longrightarrow \Lambda_t,
\]

where \texttt{Lambda\_t} is a log measure.

\subsubsection{3.3 Compatibility and allocation}\label{compatibility-and-allocation}

Processor specifications and operating constraints define which parts of
\texttt{Lambda\_t} are commercially useful. A partial-transport problem chooses an
allocation measure \texttt{Gamma\_t}.

\[
\Lambda_t \longrightarrow \Gamma_t.
\]

The allocation may leave logs unmatched, demand unmet, and operator capacity
idle. These inequalities define partial matching.

\subsubsection{3.4 Observation and learning}\label{observation-and-learning}

Observations update a posterior over physical state and structural parameters.
That posterior becomes the sufficient state for sequential decisions.

The full conceptual loop is therefore

\begin{Shaded}
\begin{Highlighting}[]
\NormalTok{latent tree state}
\NormalTok{    {-}\textgreater{} growth / mortality / recruitment}
\NormalTok{    {-}\textgreater{} harvest / bucking}
\NormalTok{    {-}\textgreater{} log distribution}
\NormalTok{    {-}\textgreater{} processor/operator matching}
\NormalTok{    {-}\textgreater{} economic outcomes and observations}
\NormalTok{    {-}\textgreater{} posterior update}
\NormalTok{    {-}\textgreater{} next decision}
\end{Highlighting}
\end{Shaded}

\subsection{4. Distinct economic objects}\label{distinct-economic-objects}

Several distinctions are essential.

\subsubsection{Biological productivity and commercial availability}\label{biological-productivity-and-commercial-availability}

Fast growth can produce logs outside the target specifications.

\subsubsection{Commercial availability and processor compatibility}\label{commercial-availability-and-processor-compatibility}

A large harvestable log volume may be incompatible with a processor's diameter,
length, species, moisture, or defect constraints.

\subsubsection{Compatibility and delivered surplus}\label{compatibility-and-delivered-surplus}

A compatible log can still be uneconomic if transport and operating costs erase
its value.

\subsubsection{Delivered surplus and negotiated forest-gate price}\label{delivered-surplus-and-negotiated-forest-gate-price}

The seller's outside option and bargaining environment can affect the realized
price even after technical matching is known.

\subsubsection{Physical coverage and information value}\label{physical-coverage-and-information-value}

The value of mapped hectares depends on their effect on a decision. Useful
observations reduce uncertainty in decision-relevant regions of state space.

\subsection{5. Structural model versus reduced-order model}\label{structural-model-versus-reduced-order-model}

The structural model is intentionally richer than current data can fully
identify. The paper later introduces a reduced-order commercial projection
based on hectares, paying accounts, matched supply, and operating cost.

The models are related by

\[
\text{structural state} \xrightarrow{\text{projection}} \text{reduced model},
\]

The reduced model is an operational approximation. A scalar ``visible tonnes''
parameter can be used until data support estimated growth, bucking, and matching
operators.

\subsection{6. Sequential control}\label{sequential-control-1}

The decision maker acts while estimates remain uncertain. Sensing, contracting,
and market exposure affect current payoffs and future information.

The action space can include:

\begin{itemize}
\tightlist
\item
  sensing or field verification;
\item
  designed experiments;
\item
  matching exposure;
\item
  contract or posted-price variables;
\item
  coverage intensity; and
\item
  discrete regime changes.
\end{itemize}

Thus the state must be defined with the eventual decision problem in mind.
There is little value in estimating a variable accurately if no feasible action
changes as a result. Conversely, a noisy variable near a processor grade
threshold can be highly decision-relevant.

\subsection{7. State, observations, and implementation}\label{state-observations-and-implementation}

The model separates \textbf{state definition}, \textbf{observation}, and \textbf{implementation}.
Individual-tree IDs support repeated measurements and remain outside the limiting
measure. Earth-observation features enter through the observation model and
become tree attributes only after calibration.

Scalar and discrete examples are explicit projections of the structural model.

\subsection{8. Mathematical scope}\label{mathematical-scope}

The model preserves the distribution of economically relevant states through the
biological, technical, and market operators. Measure-valued dynamics describe the
physical state, unbalanced transport describes feasible allocation, and Bayesian
control describes sequential decisions under partial observation.
